\section{Recursive Implementation of atoi}
\label{sec:rec-atoi}

\problemstatement{Implement (recursively) the C-style $\textit{atoi}$
function: given a string, skip any leading whitespace, then an optional
$\texttt{+}$ or $\texttt{-}$ sign, then consume as many leading digits as
possible, and return the resulting integer -- clamped to the 32-bit
signed integer range $[-2^{31}, 2^{31}-1]$ if it would otherwise
overflow. If no valid number can be parsed at all, return $0$.}

\begin{patternbox}
This problem has no clever algorithmic trick -- the difficulty is purely
in correctly handling several small edge cases (whitespace, sign,
leading zeros, early non-digit characters, overflow) using recursion
instead of an iterative loop. The keyword to notice: \emph{``recursive
implementation of''} an otherwise simple iterative task is an exercise in
expressing a left-to-right scan as a chain of recursive calls, each
processing one character and delegating the rest to a sub-call.
\end{patternbox}

\subsection*{Understanding the problem carefully}

The parsing naturally splits into three independent phases, each best
handled by its own small recursive (or simple iterative) helper: skipping
whitespace, consuming an optional sign, and then accumulating digits one
at a time. The digit-accumulation phase is the one most naturally
expressed recursively: at each step, if the current character is a
digit, fold it into a running accumulated value and recurse on the rest
of the string; the moment a non-digit (or the string's end) is reached,
stop and return the accumulated value so far.

\begin{ideabox}[An Accumulator Passed Down Through the Recursion]
$\textsc{ParseDigits}(s, i, \textit{acc})$: if $i$ is out of bounds or
$s[i]$ is not a digit, return $\textit{acc}$ (base case -- nothing more
to consume). Otherwise, compute $\textit{newAcc} = \textit{acc} \times 10
+ (s[i] - \texttt{'0'})$, clamp $\textit{newAcc}$ immediately if it has
already exceeded the representable range (to avoid undefined-behavior
overflow in the accumulator itself), and recurse:
$\textsc{ParseDigits}(s, i+1, \textit{newAcc})$.
\end{ideabox}

\subsection*{Why clamping during accumulation, not just at the end, matters}

If the accumulator is allowed to grow without bound while scanning a
string of many digits (e.g.\ a string of $50$ digits), it can overflow
even a 64-bit accumulator long before the recursion finishes, leading to
incorrect (wrapped-around) results rather than a clean clamp. Checking
and clamping the accumulator the moment it \emph{exceeds} the valid
32-bit range -- and, once clamped, simply continuing to consume (and
ignore) any remaining digits so the function still terminates at the
string's actual end of digits -- avoids this, ensuring the final
returned value is always either the true parsed integer or correctly
clamped to $\pm 2^{31}{-}1$ /$-2^{31}$.

\subsection*{Algorithm}

\begin{enumerate}
  \item Skip leading whitespace characters (advance an index $i$ while
        $s[i]$ is a space).
  \item If $s[i]$ is $\texttt{+}$ or $\texttt{-}$: record the sign,
        advance $i$.
  \item Call $\textsc{ParseDigits}(s, i, 0)$ to accumulate the digit
        value, clamping along the way as described above.
  \item Apply the recorded sign to the result, then clamp the
        signed result to $[-2^{31}, 2^{31}-1]$ one final time.
\end{enumerate}

\subsection*{Dry run}

\dryrun{Take $s = \texttt{"  -042abc"}$.}

\begin{itemize}
  \item Skip whitespace: $i$ advances past the two leading spaces.
  \item Sign: $s[i]=\texttt{'-'}$, record negative, advance $i$.
  \item $\textsc{ParseDigits}$: consume \texttt{0} ($\textit{acc}=0$),
        \texttt{4} ($\textit{acc}=4$), \texttt{2} ($\textit{acc}=42$);
        next character is \texttt{a}, not a digit -- stop, return $42$.
  \item Apply sign: $-42$. Well within range, no clamping needed.
\end{itemize}

Result: $-42$.

\subsection*{C++ implementation}

\begin{lstlisting}
#include <bits/stdc++.h>
using namespace std;

long long parseDigits(const string& s, int i, long long acc) {
    if (i >= (int)s.size() || !isdigit(s[i])) return acc;

    long long newAcc = acc * 10 + (s[i] - '0');
    if (newAcc > INT_MAX) newAcc = (long long)INT_MAX + 1; // clamp marker
    return parseDigits(s, i + 1, newAcc);
}

int myAtoi(string s) {
    int i = 0, n = s.size();
    while (i < n && s[i] == ' ') i++;

    int sign = 1;
    if (i < n && (s[i] == '+' || s[i] == '-')) {
        if (s[i] == '-') sign = -1;
        i++;
    }

    long long result = parseDigits(s, i, 0) * sign;

    if (result > INT_MAX) return INT_MAX;
    if (result < INT_MIN) return INT_MIN;
    return (int)result;
}

int main() {
    cout << myAtoi("  -042abc") << endl;          // -42
    cout << myAtoi("91283472332") << endl;        // 2147483647 (clamped)
    return 0;
}
\end{lstlisting}

\begin{complexitybox}
\textbf{Time Complexity.} $O(L)$, where $L$ is the length of the input
string -- each character is examined at most once across the
whitespace-skip, sign-check, and digit-accumulation phases.
\textbf{Space Complexity.} $O(L)$ for the digit-accumulation recursion
stack in the worst case (a string of all digits); this could be reduced
to $O(1)$ by rewriting the accumulation phase as an iterative loop, which
is how production implementations typically do it.
\end{complexitybox}

\begin{takeawaybox}
Not every recursive exercise hides a deep algorithmic insight -- some, like
this one, exist purely to practice expressing a sequential scan (skip,
check, accumulate, stop) as a chain of base-case-and-recurse calls, with
an accumulator threaded through as an explicit parameter rather than a
mutable loop variable. This accumulator-passing style is the same
technique used for tail-recursive functions in functional programming.
\end{takeawaybox}
